\begingroup
\renewcommand{\arraystretch}{1.45}
\Needspace{45\baselineskip}
\InterpolationArticle{27}
Für den Fall also, dass die Anzahl der vorgegebenen Werte der Funktion $X$, die eine vollständige Periode bilden, eine zusammengesetzte Zahl $=\pi=\mu\nu$ ist, haben wir in den Artikeln 25, 26 gelernt, durch Zerlegung jener Periode in $\nu$ Perioden von $\mu$ Gliedern dieselbe Funktion herzuleiten, die allen gegebenen Werten genügt, wie sie bei unmittelbarer Anwendung der allgemeinen Theorie auf die gesamte Periode entstehen würde. Dass die erstere Methode die Mühe der mechanischen Rechnung stärker vermindert, wird aber die Praxis denjenigen lehren, der es versucht. Es wird nun keiner weiteren Erklärung bedürfen, wie jene Zerlegung noch weiter ausgedehnt und auf den Fall angewandt werden kann, dass die Gesamtzahl aller vorgegebenen Werte aus drei oder mehr Faktoren zusammengesetzt ist, etwa wenn die Zahl $\mu$
\SourceLeaf{317}
wieder zusammengesetzt wäre; in diesem Fall lässt sich offenbar jede Periode von $\mu$ Gliedern in mehrere kleinere Perioden unterteilen. Von den vielen anderen Bemerkungen, die der praktischen Anwendung dienen, werden wir hier nur die folgenden berühren. Die Werte solcher Summen wie
\[
\begin{aligned}
&A\cos\lambda a+B\cos\lambda b+C\cos\lambda c\\[5pt]
&\qquad{}+D\cos\lambda d+\text{usw.}=p\\[5pt]
&A\sin\lambda a+B\sin\lambda b+C\sin\lambda c\\[5pt]
&\qquad{}+D\sin\lambda d+\text{usw.}=q
\end{aligned}
\]
(wobei wir wie zuvor $b-a=c-b=d\quad c\ \text{usw.}=\tfrac1\mu360^{\mathord{0}}=\Delta$ voraussetzen)
lassen sich häufig, sofern noch nicht $a=0$ gilt, bequemer durch die Formeln
\[
\begin{aligned}
&A+B\cos\lambda\Delta+C\cos2\lambda\Delta\\[5pt]
&\qquad{}+D\cos3\lambda\Delta+\text{usw.}=P\\[5pt]
&B\sin\lambda\Delta+C\sin2\lambda\Delta\\[5pt]
&\qquad{}+D\sin3\lambda\Delta+\text{usw.}=Q
\end{aligned}
\]
$p=P\cos a-Q\sin a$, $q=P\sin a+Q\cos a$ bestimmen, oder auch, indem man
$\frac QP=\operatorname{tang}\varphi$ und $\frac P{\cos\varphi}=\frac Q{\sin\varphi}=R$ setzt,
durch die Formeln $p=R\cos(\varphi+a)$, $q=R\sin(\varphi+a)$.

\endgroup
